% Modernised reading — fonds Alexandre Grothendieck, shelfmark 29, batch 8
% (pages 141-160).
%
% This is an INTERPRETATION, not a transcription. It re-reads the pages in
% today's notation and vocabulary, states the hypotheses the manuscript leaves
% implicit, and drops the critical apparatus so that the mathematics reads
% straight through. Everything the brackets and underlines carried moves into a
% footnote. In any dispute of substance the transcription governs: whoever
% cites Grothendieck must cite batch-08.fr.tex.
%
% Scope: the folder was read WHOLE before any of this was written — all eleven
% batches, pages 1-216, in one pass — and one file was then written per batch.
% Two things this file could not have said batch by batch: that the axioms of
% pages 157-158 and those of page 209 are two states of ONE list, seventy pages
% apart and in different numberings; and that the word « connexe » replacing
% « indécomposable » interlinearly on page 157 is the word that pages 190-216
% then use throughout.
%
% This batch contains the faintest writing in the folder (pages 141, 155, 156).
% Where the page does not read, this reading says so and does not fill in.
%
% Departures from the page, each footnoted where it occurs:
%   - page 154 draws the residue-field tower k(x'), k(x), k(y) top-down, the
%     reverse of the inclusions its own argument uses; it is drawn here in the
%     order the argument requires;
%   - page 142 writes « les s-gpes G de C d'un G de C », the same letter twice
%     in one sentence; the two are separated here;
%   - Gal 5 is stated on the page « sous les conditions de Gal 3 » where the
%     margin's own remark shows what is meant.
%
% Pass: Opus 5 (claude-opus-5), 2026-08-29 — first pass, unchecked against the
% pages by a human.
\documentclass[11pt,a4paper]{article}
% Le préambule commun à toutes les transcriptions du fonds Grothendieck.
%
% Il tient en une idée : **l'incertitude se compose, elle ne se résout pas.**
% Une transcription qui lisse un mot douteux a détruit la seule chose pour
% laquelle on la faisait — un lecteur doit pouvoir savoir, à chaque endroit,
% ce qui était sur le feuillet et ce qui a été ajouté. Les six macros qui
% suivent sont donc l'appareil critique, et non de la décoration.
%
% Les mêmes commandes sont reconnues par `scripts/render.mjs`, qui produit la
% vue de lecture du site. Une macro ajoutée ici doit l'être aussi là-bas,
% faute de quoi le rendu échouera bruyamment — ce qui est voulu : un rendu
% silencieusement faux est pire qu'un rendu absent.

\ProvidesPackage{grothendieck}[2026/08/08 Transcriptions du fonds Grothendieck]

\usepackage{amsmath,amssymb,amsthm}
\usepackage{mathrsfs}
\usepackage{stmaryrd}
% Les diagrammes commutatifs sont partout dans ce fonds : c'est la langue
% même de la géométrie algébrique de Grothendieck, et une transcription qui
% les rendrait en prose ne transcrirait rien.
\usepackage{tikz-cd}
% Français par défaut — c'est la langue des feuillets — mais l'anglais est
% chargé aussi : la traduction et le résumé sont des documents anglais, et la
% typographie française (espaces avant les deux-points) y serait une faute.
% Ces documents basculent par \selectlanguage{english} en tête de corps.
\usepackage[english,french]{babel}
\usepackage{fontspec}
\usepackage{geometry}
\geometry{a4paper,margin=25mm,marginparwidth=28mm}
\usepackage{marginnote}
\usepackage{xcolor}
% ulem, not soul. `soul`'s \st and \ul are built on a character-by-character
% scan that XeLaTeX's font machinery breaks: the document fails with an
% undefined control sequence rather than degrading. ulem does the same two jobs
% and is Unicode-engine safe. `normalem` stops it hijacking \emph, which the
% transcriptions use for Grothendieck's own emphasis.
\usepackage[normalem]{ulem}
\usepackage{hyperref}
\hypersetup{colorlinks=true,linkcolor=black,urlcolor=black,pdfborder={0 0 0}}

% --- Métadonnées du lot -----------------------------------------------------
% Reprises telles quelles par le script de rendu, qui les affiche en tête de
% la vue de lecture. Elles fixent ce que le fichier prétend couvrir : sans
% elles, une transcription flotte sans point d'attache dans un fonds de seize
% mille feuillets.

\newcommand{\folder}[1]{\gdef\grothFolder{#1}}
\newcommand{\batch}[1]{\gdef\grothBatch{#1}}
\newcommand{\leaves}[2]{\gdef\grothFirst{#1}\gdef\grothLast{#2}}
\newcommand{\foldertitle}[1]{\gdef\grothFoldertitle{#1}}
\gdef\grothFolder{?}\gdef\grothBatch{?}\gdef\grothFirst{?}\gdef\grothLast{?}\gdef\grothFoldertitle{}

% --- Le résumé --------------------------------------------------------------
%
% Il ouvre la lecture modernisée, sous le titre « Résumé », et il est composé
% en petit. Ce n'est pas un ornement : c'est le seul passage du document qui
% ne soit pas des mathématiques, et un lecteur doit pouvoir le reconnaître
% avant de l'avoir lu — puis le sauter s'il connaît déjà le sujet, ou s'y
% arrêter s'il ne le connaît pas. Le filet qui le suit marque la reprise du
% corps.

\newenvironment{resume}
  {\small\setlength{\parskip}{0.45em}}
  {\par\medskip\hrule\medskip}

% --- Mots-clés ---------------------------------------------------------------
%
% En anglais, en clôture du résumé de la lecture modernisée : le vocabulaire
% moderne sous lequel chercher ce que les feuillets construisent. Le manifeste
% du site (`npm run manifest`) les extrait pour étiqueter la cote entière —
% cette ligne est la source unique des tags, il n'y a pas de fichier à côté.
\newcommand{\keywords}[1]{\par\smallskip\noindent
  {\footnotesize\textsc{Keywords} --- \textit{#1}}\par}

% --- L'appareil critique ----------------------------------------------------

% Le numéro de feuillet, en marge. C'est l'ancre qui permet de régler un
% désaccord sur une formule en regardant *un* feuillet plutôt qu'une chemise
% de six cents. Le site s'en sert aussi pour tourner les pages du fac-similé
% quand on fait défiler la transcription.
\newcommand{\leaf}[1]{%
  \marginnote{\footnotesize\color{gray!70}\textsf{#1}}%
  \label{leaf:#1}\ignorespaces}

% Illisible : on ne devine pas. Un mot inventé qui se lit comme les autres est
% le pire résultat possible.
\newcommand{\ill}{\textcolor{red!60!black}{[\,\ldots\,]}}

% Lecture douteuse : le mot est donné, et signalé comme incertain.
\newcommand{\uncertain}[1]{\uline{#1}}

% Ajout de l'éditeur — une lettre manquante, un mot sous-entendu.
\newcommand{\add}[1]{\textcolor{blue!50!black}{[#1]}}

% Biffé par Grothendieck lui-même. Ce qu'il a rayé fait partie du document :
% une rature dit souvent où la pensée a changé de direction.
\newcommand{\struck}[1]{\sout{#1}}

% Note du transcripteur, sur l'état matériel du feuillet.
\newcommand{\note}[1]{\par\smallskip{\small\color{gray!60!black}\textit{[#1]}}\par\smallskip}

% Note marginale de Grothendieck, distincte du corps du texte.
\newcommand{\marginal}[1]{\par\smallskip\noindent\hspace{1em}%
  {\small\color{gray!60!black}#1}\par\smallskip}

% --- Titre ------------------------------------------------------------------

\AtBeginDocument{%
  \begin{center}
    {\large\bfseries Fonds Alexandre Grothendieck --- cote n\textsuperscript{o}~\grothFolder}\\[2pt]
    {\itshape\grothFoldertitle}\\[4pt]
    {\small feuillets \grothFirst--\grothLast\ (lot \grothBatch)}\\[2pt]
    {\footnotesize\color{gray!60!black}Transcription établie d'après le fac-similé numérisé
     par l'Université de Montpellier.\\
     Les lectures incertaines sont soulignées, les illisibles notées [\,\ldots\,],
     les ajouts entre crochets.}
  \end{center}
  \vspace{1em}}

\endinput


\folder{29}
\batch{8}
\pages{141}{160}
\dating{1959-1969}
\watermark{Édition de démonstration\\interprétation personnelle de l'œuvre}
\foldertitle{Groupe fondamental [Autour de SGA 4 et SGA 7] — lecture modernisée}

\begin{document}

\section*{Résumé}

\begin{resume}
La chemise porte « Compléments SGA / 1960 », et le lot fait ce que ce titre
promet : il achève ce que la lettre de 1959 avait laissé en suspens, puis il
écrit, pour la première fois dans ce carton, la liste d'axiomes qui définit une
\emph{catégorie galoisienne}.

Cette liste est ce qu'il faut regarder. Toute la théorie de Galois classique
tient dans une observation : pour connaître un corps, il n'est pas nécessaire
d'en connaître les éléments, il suffit de connaître la collection de ses
extensions finies, avec la façon dont elles s'emboîtent. Grothendieck demande
alors : \emph{quelles propriétés une collection d'objets doit-elle avoir pour
qu'elle soit la collection des revêtements de quelque chose} ? Six axiomes
suffisent. Il faut que les sommes et les produits existent ; que tout objet se
décompose de façon essentiellement unique en morceaux indécomposables --- il
raye ce mot partout et écrit \emph{connexes} au-dessus, et c'est le mot qui
restera ; que les morceaux d'un objet qui s'envoient isomorphiquement sur un
objet connexe correspondent exactement aux sections ; qu'un foncteur vers les
ensembles finis, le \emph{foncteur fibre}, respecte tout cela et ne déclare
isomorphe que ce qui l'est ; et que les objets galoisiens admettent des
quotients. Sous ces conditions, la collection est celle des ensembles finis sur
lesquels opère un groupe profini, et ce groupe est le groupe fondamental.

Ce qui rend ces pages intéressantes à lire est qu'on y voit la liste se faire.
Un axiome est écrit puis biffé en entier ; les numéros sont réattribués en
cours de route ; et une seconde version est tentée aussitôt après, sur une
catégorie munie non d'un foncteur fibre mais d'une sous-catégorie de
\emph{points} et d'une classe de morphismes « non ramifiés » --- c'est-à-dire
en remplaçant le foncteur par la géométrie. Le carton reprendra la question
cinquante pages plus loin, page 209, avec une autre liste et une autre
numérotation, pour une classe d'objets plus large.

Le lot contient aussi deux résultats fins. Le premier, dû à Serre et crédité
comme tel sur la page, est la clef de toute la théorie : un revêtement fini
étale de rang $n$ est toujours de la forme $P/H$ pour un revêtement galoisien
$P$ de groupe le groupe symétrique et $H$ le stabilisateur d'un point --- on
l'obtient en considérant les \emph{$n$-uplets de points distincts} au-dessus
d'un même point de la base. Le second est un théorème de Künneth : le groupe
fondamental d'un produit fibré s'envoie toujours \emph{surjectivement} sur le
produit fibré des groupes fondamentaux, et il faut une condition
supplémentaire, explicitée ici, pour que ce soit un isomorphisme.

Les noms modernes à chercher sont \emph{catégorie galoisienne}, \emph{foncteur
fibre}, \emph{morphisme fini étale}, \emph{formule de Künneth pour le groupe
fondamental}.

\keywords{Galois category, axiomatic Galois theory, finite etale morphism, quotient by a finite group, Kunneth formula for the fundamental group}
\end{resume}

\pagerange{141}{142}
\section*{Le pro-objet $\pi_{1}^{\mathcal{C}}(S,\xi)$ (pages 141 et 142)}

\subsection*{La fin de la proposition 4}

La démonstration commencée page 140 s'achève par extension de la base. On se
donne $(S_{1}, a_{1}) \to (S, a)$ ; la situation devient celle de $G_{1}$,
$G'_{1}$, $(P_{1}, \xi_{1})$, $(Q_{1}, \eta_{1})$, et le fibré $P_{1}$
s'identifie alors à $G'_{1}$ muni de son marquage $\varepsilon'_{1}$, tandis que
$Q_{1}$ s'identifie à $G'_{1} \times_{S_{1}} G'_{1}$ muni de
$(\varepsilon'_{1}, \varepsilon'_{1})$. Il suffit alors de voir que
$(P'_{1}, \xi'_{1}) \to (Q_{1}, \eta_{1})$ est l'inclusion diagonale
$G'_{1} \to G'_{1} \times_{S_{1}} G'_{1}$.

Notant $w$ l'isomorphisme $(P, \xi) \times_{u} G' \to (P, \xi) \times_{v} G'$
et $w_{1}$ son extension à $S_{1}$, l'assertion revient à dire que $w_{1}$
transporte la section naturelle de $P_{1} \times_{u_{1}} G'_{1}$ sur celle de
$P_{1} \times_{v_{1}} G'_{1}$ ; et cela résulte de l'unicité de
$P \to P \times_{u} G'$ compatible au marquage.

\begin{tikzcd}[column sep=large]
P \arrow[r, "\mathrm{can}"] \arrow[rr, bend right=30, "\mathrm{can}"'] & P \times_u G' \arrow[r, "w"] & P \times_v G'
\end{tikzcd}

\subsection*{Le pro-objet}

Restent deux propriétés formelles. D'abord la multiplicativité,
\[
  \mathfrak{Z}(G \times G') \;=\; \mathfrak{Z}(G) \times \mathfrak{Z}(G') ,
\]
et ensuite la condition minimale sur les sous-groupes appartenant à
$\mathcal{C}$ d'un objet de $\mathcal{C}$.\footnote{La phrase de la page ---
« Les s-gpes $G \in \mathcal{C}$ d'un $G \in \mathcal{C}$ satisfaisant la
condition minimale » --- fait servir la même lettre deux fois, pour le
sous-groupe et pour l'objet. La transcription le signale et le conserve ; ici
les deux sont séparés, et il s'agit bien de la condition (iii) de la lettre à
Serre : la chaîne descendante des sous-groupes est stationnaire.}

\emph{Ces propriétés formelles suffisent} pour construire un pro-objet de
$\mathcal{C}$, au moyen d'un système projectif d'épimorphismes, que l'on note
\[
  \pi_{1}^{\mathcal{C}}(S, \xi) .
\]

C'est le même objet que la lettre de 1959 construit page 126, obtenu ici sans
passer par les couples minimaux, mais par les seules propriétés du foncteur.
Deux exemples sont amorcés et laissés inachevés : $\mathcal{C}$ une classe de
groupes définis par un groupe abstrait --- « on retrouve le groupe
fondamental » --- et $\mathcal{C}$ provenant d'un corps
$k$.\footnote{Les deux pages 141 et 142 sont au crayon très pâle et sont les
plus difficiles du carton ; la transcription y laisse beaucoup d'illisible et
cette lecture ne le comble pas. Ce qui est restitué ici sont les énoncés
portants, non les phrases qui les joignent.}

\pagerange{144}{146}
\section*{Données de descente : la relation de cocycle (pages 144 à 146)}

Deux feuillets écrivent la condition de cocycle pour une donnée de descente le
long d'un morphisme, dans le langage des groupoïdes fondamentaux.

On dispose d'un $S' \to S$ et de ses puissances fibrées, avec les projections
$p_{i}^{\alpha}$ de $\pi''_{\alpha}$ dans $\pi'$ et $p_{ij}^{\alpha\beta}$ de
$\pi'''_{\beta}$ dans $\pi''_{\alpha}$, indexées par les trois couples
$(1,2)$, $(2,3)$, $(1,3)$. Les relations entre elles font intervenir des
automorphismes intérieurs, et le carré qui les résume est celui-ci :

\begin{tikzcd}[column sep=large]
E \arrow[r, "A_{\alpha'}"] \arrow[d, "a'_{\beta}"'] & E \arrow[r, "a''_{\beta}", "\sim"'] & E \arrow[d, "A_{\alpha''}"] \\
E \arrow[r, "A_{\alpha'''}"'] & E \arrow[r, "a'''_{\beta}"', "\sim"] & E
\end{tikzcd}

d'où la relation, que la page encadre :
\[
  \boxed{\; A_{\alpha''}\, a''_{\beta}\, A_{\alpha'} \;=\;
  a'''_{\beta}\, A_{\alpha'''}\, a'_{\beta} \;}
\]

C'est la condition de cocycle, écrite comme la commutation d'un rectangle. La
page 146 la reprend en termes de \emph{chemins} --- $g_{s''} = g_{s'''}g_{s'}$,
d'où $g = 1$ pour des chemins triviaux --- puis abandonne : elle est barrée
presque entièrement. Ce qui reste au bas de la feuille est le diagramme
simplicial $S \leftarrow S' \leftarrow S'' \leftarrow S'''$ avec les
$R_{i}$ au-dessus, et trois mots isolés.

Cette relation reviendra sous une forme réglée à la page 180, où les données de
descente sur un faisceau constant sont décrites comme des applications
$f : K_{1} \to G$ vérifiant $f \circ p_{1} = (f \circ p_{0})(f \circ p_{2})$.
C'est la même identité, et c'est là seulement qu'elle est écrite proprement.

\pagerange{148}{150}
\section*{« Compléments aux notes IHES » : le $\pi_{1}$ d'un produit fibré (pages 148 à 150)}

\subsection*{L'énoncé}

Soient $X$ et $Y$ deux schémas sur $S$, avec $X$ propre, séparable et surjectif
sur $S$, et $Y$ non vide. Soit $c$ un point géométrique de $X \times_{S} Y$,
d'images $a$, $b$, $d$ dans $X$, $Y$, $S$. On dispose de l'homomorphisme
canonique
\[
  \pi_{1}(X \times_{S} Y,\, c) \longrightarrow
  \pi_{1}(X, a) \times_{\pi_{1}(S, d)} \pi_{1}(Y, b) .
\]

\textbf{Il est toujours surjectif.} Il est \emph{bijectif} dès que la condition
suivante est vérifiée : il existe un $S'$ non vide, étale sur $S$, tel que
$Y' = Y \times_{S} S'$ admette une section sur $S'$ --- c'est-à-dire, ce qui
revient au même, tel qu'il existe un $S$-morphisme $S' \to Y$ --- et tel que
$\pi_{1}(X') \to \pi_{1}(X)$ soit injectif.

Le cas particulier qui sert : lorsque $S = \operatorname{Spec} k$, on obtient
\[
  \pi_{1}(X \times_{k} Y) \;\simeq\; \pi_{1}(X) \times_{\pi_{1}(k)}
  \pi_{1}(Y)
\]
pour $X$ propre universellement connexe sur $k$ et $Y$ simplement connexe sur
$k$, sans plus.

Il faut mesurer ce que la condition coûte. Elle demande un point rationnel de
$Y$, au moins après une extension étale de la base --- ce qui n'est pas
automatique --- et l'injectivité de $\pi_{1}(X') \to \pi_{1}(X)$. Sans elle,
seule la surjectivité subsiste, et la différence est réelle : c'est elle qui
sépare le théorème de Künneth pour $\pi_{1}$ de la simple existence d'un
homomorphisme.

\subsection*{La démonstration}

Elle est page 150, et c'est une échelle de suites exactes. Les deux noyaux
\[
  \operatorname{Ker}\bigl[\pi_{1}(X) \to \pi_{1}(S)\bigr]
  \qquad\text{et}\qquad
  \operatorname{Ker}\bigl[\pi_{1}(X \times_{S} Y) \to \pi_{1}(Y)\bigr]
\]
s'identifient tous deux à l'image de $\pi_{1}(\bar{X}_{s})$ --- le groupe
fondamental de la fibre géométrique --- et la comparaison des deux lignes donne
le résultat.

\begin{tikzcd}[column sep=small, row sep=small, nodes={font=\scriptsize}]
& \operatorname{Im}(\pi_1(\bar{X}_s) \to \pi_1(X)) \arrow[d, no head] & & & \\
e \arrow[r] & \operatorname{Ker}[\pi_1(X) \to \pi_1(S)] \arrow[r, "\sim"] & \pi_1(X) \times_{\pi_1(S)} \pi_1(Y) \arrow[r] & \pi_1(Y) \arrow[r] & e \\
e \arrow[r] & \operatorname{Ker}[\pi_1(X \times_S Y) \to \pi_1(Y)] \arrow[u] \arrow[d, no head] \arrow[r] & \pi_1(X \times_S Y) \arrow[u] \arrow[r] & \pi_1(Y) \arrow[u] \arrow[r] & e \\
& \operatorname{Im}(\pi_1(\bar{X}_s) \to \pi_1(X \times_S Y)) & & &
\end{tikzcd}

C'est la propreté de $X$ sur $S$ qui fait marcher les deux identifications :
elle donne la suite exacte d'homotopie, dont les deux lignes sont deux
instances.

La page 149, entre les deux, ne porte qu'un croquis : $X \to S$ avec $S'$ à
côté. C'est le changement de base dont les deux autres pages discutent, et
c'est tout ce que la feuille contient.

\pagerange{151}{151}
\section*{Le topos des $\pi$-ensembles (page 151)}

Une page entière barrée de trois traits diagonaux, et elle mérite d'être lue
malgré la biffure : elle vérifie que la catégorie des ensembles finis à
opérateurs est un topos et que sa cohomologie est celle du groupe.

Soit $\pi$ un groupe et $\mathcal{E}^{\pi}$ la catégorie des ensembles finis où
$\pi$ opère à gauche, l'objet final étant le point $e$, l'objet initial
$\varnothing$. Le $\pi$-ensemble $\pi_{s} = \pi$ (translations à droite) a pour
groupe d'automorphismes $\pi$ lui-même, par $x \mapsto xg^{-1}$.

Un faisceau d'ensembles sur $\mathcal{E}^{\pi}$ est alors la même chose qu'un
ensemble fini $E$ où $\pi$ opère, via
\[
  E \;=\; F(\pi_{s}) , \qquad
  F(U) \;=\; \operatorname{Hom}_{\pi}(U, E) .
\]
Plus précisément, si $U = \coprod_{i} \pi/\pi_{i}$ est la décomposition en
orbites, alors la catégorie fibrée $\mathcal{F}^{\pi}$ vérifie
$\mathcal{F}_{U} \simeq \prod_{i} E^{\pi_{i}}$ ; et
\[
  H^{0}(F) = \Gamma(F) = E^{\pi} , \qquad
  H^{1}(F) \simeq H^{1}(\pi, E) , \qquad
  H^{i}(F) \simeq H^{i}(\pi, E) ,
\]
la dernière ligne valant pour $F$ et $E$ commutatifs.

C'est le \emph{topos classifiant} $B\pi$, et l'identification de sa cohomologie
à celle du groupe. La page 209 s'en servira sans le redémontrer : « prenant
pour $\mathcal{T}$ le topos $(B\pi_{i})$ ».

\pagerange{152}{156}
\section*{Un morphisme fini étale de rang $n$ est un $P/H$ (pages 152 à 156)}

\subsection*{La proposition 3.2}

\textbf{Proposition.} Soit $f : X \to Y$ un morphisme fini. Les conditions
suivantes sont équivalentes.

\begin{enumerate}
  \item[(i)] $f$ est étale, et $X$ est de rang localement égal à $n$ sur $Y$.
  \item[(ii)] Il existe $Y_{1} \to Y$ fidèlement plat tel que
    $X \times_{Y} Y_{1}$ soit $Y_{1}$-isomorphe à $Y_{1} \times I_{n}$, où
    $I_{n}$ est un ensemble à $n$ éléments.
  \item[(ii bis)] La même chose avec $Y_{1}$ un revêtement galoisien, de groupe
    le groupe symétrique $\mathfrak{S}_{n}$.
  \item[(iii)] $X$ est isomorphe à un $P/H$, où $P$ est un revêtement étale
    galoisien de $Y$ de groupe $G$ et $H$ un sous-groupe de $G$.
  \item[(iii bis)] La même chose avec $G = \mathfrak{S}_{n}$ et $H$ le
    stabilisateur d'un point de $I_{n}$.
\end{enumerate}

Entre (ii) et (ii bis), et entre (iii) et (iii bis), la page écrit
« kif -- kif » : c'est la même chose.

\subsection*{La démonstration de Serre}

L'implication (i) $\Rightarrow$ (iii bis) est créditée sur la page : « la
démonstration qui suit [est] due à J.-P.\ Serre ». Elle vaut d'être écrite en
entier, car elle est le mécanisme de toute la théorie.

Soit $X^{n} = X \times_{Y} \cdots \times_{Y} X$ le produit fibré de $n$
copies, sur lequel $\mathfrak{S}_{n}$ opère par permutation des facteurs. Soit
$P \subset X^{n}$ le complémentaire du lieu où deux coordonnées coïncident ---
c'est-à-dire l'ensemble des $n$-uplets de points \emph{deux à deux distincts}
au-dessus d'un même point de $Y$. C'est une partie \emph{ouverte et fermée} de
$X^{n}$, parce que le lieu de coïncidence en est une (la diagonale d'un
morphisme étale est ouverte et fermée) ; elle est stable sous
$\mathfrak{S}_{n}$, qui y opère \emph{sans point fixe} par construction. Enfin,
$X$ étant de rang $n$ sur $Y$, la fibre géométrique de $P$ au-dessus d'un point
de $Y$ est exactement l'ensemble des $n!$ bijections de $I_{n}$ sur la fibre de
$X$.

Donc $P \to Y$ est un revêtement étale galoisien de groupe
$\mathfrak{S}_{n}$. Et la dernière projection $X^{n} \to X$ est invariante sous
le stabilisateur $H = \mathfrak{S}_{n-1}$ du dernier indice, d'où
$X \simeq P/H$.\footnote{C'est la construction dite du \emph{revêtement des
repères} : $P$ est le schéma des isomorphismes $I_{n} \to X$ au-dessus de $Y$,
et le fait qu'il soit galoisien de groupe $\mathfrak{S}_{n}$ est ce qui permet
de ramener tout revêtement fini étale à un revêtement galoisien. C'est l'énoncé
qui rend la théorie de Galois abstraite non vide : sans lui, l'existence d'un
revêtement galoisien majorant un revêtement donné serait à démontrer cas par
cas.}

Les implications restantes se ferment : (iii bis) $\Rightarrow$ (iii) est
triviale ; (iii) $\Rightarrow$ (i) résulte d'un lemme antérieur ; (iii)
$\Rightarrow$ (ii bis) en prenant $Y_{1} = P$ et en utilisant
$(X \times G)/H \simeq X \times G/H$ ; et (ii bis) $\Rightarrow$ (i) est
immédiat.

\subsection*{La proposition 3.1 : l'existence de $X/G$}

\textbf{Proposition.} Soit $X$ étale sur $Y$ et quasi-compact, et $G$ un groupe
fini opérant par $Y$-automorphismes. Alors $X/G$ existe, $X \to X/G$ est fini,
et $X/G \to Y$ est \emph{étale}.

\emph{Démonstration.} Soit $I$ l'ensemble des composantes connexes de $X$, sur
lequel $G$ opère ; en regroupant les composantes par orbite, $X$ est la somme
des $X_{(\alpha)} = \bigcup_{i \in \alpha} X_{i}$, stables sous $G$, et
$X/G = \bigsqcup_{\alpha} X_{(\alpha)}/G$. On se ramène donc au cas où $G$
opère transitivement sur les composantes connexes. Soit $X_{0}$ l'une d'elles
et $G_{0}$ son stabilisateur : alors $X \simeq X_{0} \times_{G_{0}} G$ et
$X/G \simeq X_{0}/G_{0}$.

On peut alors supposer que $G_{0}$ opère fidèlement ; mais il opère aussi
\emph{sans inertie}, et il reste à voir que $X/G$ est étale sur $Y$. Faisant
une extension fidèlement plate de la base, on se ramène au cas où l'extension
résiduelle $k(x)/k(y)$ est triviale, et l'on conclut sur les complétés :
\[
  \hat{\mathcal{O}}_{y} \longrightarrow \hat{\mathcal{O}}_{x'}
  \xrightarrow{\ \sim\ } \hat{\mathcal{O}}_{x} .
\]

\begin{tikzcd}[column sep=small, row sep=small]
X \arrow[dr] & \\
& X' = X/G \arrow[d] \\
& Y \ni y
\end{tikzcd}
\qquad
\begin{tikzcd}[row sep=small]
k(x) \arrow[d, no head] \\
k(x') \arrow[d, no head] \\
k(y)
\end{tikzcd}

La tour des corps résiduels est dessinée ici dans l'ordre des inclusions,
$k(y) \subset k(x') \subset k(x)$.\footnote{La page la dessine dans l'ordre
inverse --- $k(x')$ en haut, puis $k(x)$, puis $k(y)$ --- ce que la
transcription signale. C'est un lapsus de tracé, et non une convention : la
suite de l'argument, qui passe de $\hat{\mathcal{O}}_{y}$ à
$\hat{\mathcal{O}}_{x'}$ puis à $\hat{\mathcal{O}}_{x}$, impose l'ordre écrit
ici.}

\textbf{Corollaire.} Si $X$ est étale sur $Y$ en $x$ et si $g : X \to X'$ est
un $Y$-morphisme étale en $x$, alors $X'$ est étale sur $Y$ en
$g(x)$.\footnote{Les pages 155 et 156, au crayon sur une transparence de
verso, sont les plus difficiles du lot après la 141 ; la transcription y laisse
en particulier la condition (i) de la proposition 3.1 largement illisible. Ce
qui est énoncé ci-dessus est ce que les phrases portantes autorisent, et
l'hypothèse « quasi-affine » que la page ajoute en cours de démonstration est
celle sous laquelle le quotient existe.}

\pagerange{157}{158}
\section*{« Catégorie galoisienne » : Gal 1 à Gal 6 (pages 157 et 158)}

Ce sont les deux pages les plus citables du lot. Elles écrivent la liste
d'axiomes, et l'on y voit la liste se faire : un axiome est écrit puis biffé en
entier, les numéros sont réattribués, et le mot « indécomposable » est remplacé
interlinéairement par « connexe » à chacune de ses occurrences.

\subsection*{Les données}

\begin{enumerate}
  \item[(i)] Une catégorie $\mathcal{C}$.
  \item[(ii)] Un foncteur $F : \mathcal{C} \to (\text{ensembles finis})$.
\end{enumerate}

\subsection*{Les axiomes}

\begin{itemize}
  \item[\textbf{Gal 1}] Dans $\mathcal{C}$, les sommes et les produits finis
    existent --- d'où l'existence de l'objet initial $0$ et de l'objet final
    $1$.

  \item[\textbf{Gal 4}] $F$ commute aux sommes et aux produits.
\end{itemize}

On dit alors qu'un objet $X$ est \emph{connexe} si $X = X' \sqcup X''$ implique
$X' \simeq 0$ ou $X'' \simeq 0$.

\begin{itemize}
  \item[\textbf{Gal 2}] Tout $X$ s'écrit $X \simeq \coprod_{i \in I} X_{i}$
    avec les $X_{i}$ connexes non nuls, et toute décomposition
    $X \simeq X' \sqcup X''$ provient d'une partition $I = I' \sqcup I''$. Les
    $X_{i}$, déterminés de façon essentiellement unique avec leurs flèches vers
    $X$, sont les \emph{composantes connexes} de $X$. On a $X = 0$ si et
    seulement si $I = \varnothing$.

  \item[\textbf{Gal 3}] Soit $f : X \to Y$ avec $Y$ connexe non nul. Alors les
    composantes $X_{i}$ de $X$ telles que $f$ induise un \emph{isomorphisme}
    $f_{i} : X_{i} \xrightarrow{\ \sim\ } Y$ sont en correspondance
    biunivoque avec les \emph{sections} $g : Y \to X$ de $f$, la section
    associée à $X_{i}$ étant le composé
    $Y \xrightarrow{f_{i}^{-1}} X_{i} \hookrightarrow X$. Il suffit de le
    vérifier pour $X = Z \times Y$ et $f = p_{2}$.

  \item[\textbf{Gal 5}] Si $f : X \to Y$ induit une bijection
    $F(X) \xrightarrow{\ \sim\ } F(Y)$, c'est un isomorphisme. Cela implique en
    particulier que $F(X) = \varnothing$ entraîne $X = 0$. Il suffit de le
    vérifier quand $X$ est une composante connexe d'un produit $Z \times Y$ et
    $f$ induit par $p_{2}$.

  \item[\textbf{Gal 6}] Soit $X$ connexe non nul et \emph{galoisien},
    c'est-à-dire tel que $G = \operatorname{Aut}(X)$ opère
    \emph{transitivement} dans $F(X)$. Alors pour tout sous-groupe $H$ de $G$,
    le quotient $X/H$ existe dans $\mathcal{C}$ et
    \[
      F(X/H) \;\simeq\; F(X)/H .
    \]
\end{itemize}

\subsection*{Ce qu'on voit dans la fabrication}

Un axiome antérieur --- « $F(X) = \varnothing \Rightarrow X$ est une unité à
gauche » --- est biffé en entier, et pour cause : il est \emph{conséquence} de
Gal 5, comme la page le note aussitôt entre crochets. C'est un axiome retiré
parce que redondant, et la même chose se reproduira page 211 avec l'axiome
\textbf{G5}, dont il écrira qu'il le soupçonne redondant sans le
démontrer.\footnote{Cette liste n'est pas celle de SGA 1, exposé V, § 4, dont
la numérotation et le contenu diffèrent --- là, G3 est la factorisation d'un
morphisme en un épimorphisme suivi d'un isomorphisme sur un facteur direct, et
G4 à G6 portent sur l'exactitude et la conservativité de $F$. C'est un état
antérieur, où la décomposition en composantes connexes est encore un axiome
(Gal 2) plutôt qu'une conséquence. La page 209 du même carton donnera un
troisième état, sous le nom de G0 à G5, et pour une classe d'objets plus large.}

Le remplacement de « indécomposable » par « connexe » n'est pas une correction
de style : c'est le moment où le vocabulaire de la géométrie passe dans une
axiomatique qui n'en a pas besoin, et c'est le mot que les pages 190 à 216
emploieront sans plus jamais l'expliquer.

\subsection*{La seconde version : des points au lieu d'un foncteur}

La page 158 recommence, avec un dispositif différent : une catégorie
$\mathcal{C}$, une sous-catégorie \emph{pleine} $\mathcal{P}$ dont les objets
s'appellent les \emph{points}, et une classe de morphismes dits \emph{non
ramifiés}. Les neuf conditions, numérotées $1^{\mathrm{o}}$ à
$9^{\mathrm{o}}$ dans un ordre qui n'est pas celui de leur écriture, disent :
que les sommes finies et les produits fibrés existent ; que les morphismes non
ramifiés sont stables par composition et changement de base ; qu'il y a des
composantes connexes ; que les points sont connexes, de sorte que
$\operatorname{Hom}(\xi, X \sqcup Y) = \operatorname{Hom}(\xi, X) \sqcup
\operatorname{Hom}(\xi, Y)$ ; qu'une section d'un morphisme non ramifié sur un
objet connexe est un isomorphisme sur une composante connexe ; et qu'un
morphisme non ramifié dont les fibres $\operatorname{Hom}(\xi, X) \to
\operatorname{Hom}(\xi, Y)$ sont réduites à un élément est un isomorphisme.
Enfin, la neuvième donne l'existence de $X/H$ comme en Gal 6.

Sous ces conditions on définit $\pi_{1}(X, x)$ pour $x \in
\operatorname{Hom}(\xi, X)$, et --- c'est la dernière ligne de la page ---
« c'est un foncteur \emph{covariant} en $(X, x)$ » :

\begin{tikzcd}
\xi \arrow[r, "x"] \arrow[d] & X \arrow[d] \\
\xi' \arrow[r, "x'"'] & X'
\end{tikzcd}

Le foncteur fibre a été remplacé par la famille des points ; c'est le passage
du groupe fondamental au \emph{groupoïde} fondamental, et c'est la structure
que les catégories multigaloisiennes de la page 208 axiomatiseront.

\pagerange{159}{160}
\section*{« Groupe fondamental dans les catégories » (pages 159 et 160)}

Deux dernières pages poussent l'abstraction d'un cran, et c'est le cadre le
plus large du carton.

Soit $\mathcal{C}$ une catégorie munie d'une famille $(D)$ de morphismes
satisfaisant :

\begin{enumerate}
  \item[(i)] $(D)$ contient les identités et est stable par composition ;
  \item[(ii)] $(D)$ est stable par changement de base ;
  \item[(iii)] les morphismes de $(D)$ sont des morphismes de descente
    \emph{universels stricts} ;
  \item[(iv)] les morphismes qui « descendent » par un $g \in (D)$ le font pour
    un $f \in (D)$ --- mais \emph{on ne suppose pas} qu'on puisse faire
    descendre tout morphisme par un $f \in (D)$ ;
  \item[(v)] la projection $X \times I \to X$ est dans $(D)$.
\end{enumerate}

\emph{Exemple.} $\mathcal{C}$ la catégorie des préschémas, $(D)$ les morphismes
\emph{plats et propres} --- ou plats et quasi-propres.

On développe alors la théorie des fibrés principaux dans ce cadre : les groupes
structuraux sont les $X \times G$ pour $G$ un groupe abstrait fini, et les
fibrés principaux sur $X$ de groupe $G$ sont ceux dont la projection
$P \to X$ appartient à $(D)$.

Il faut voir ce que la clause (iv) refuse. Elle dit qu'on ne demande pas à
$(D)$ d'être une classe de recouvrements au sens d'une topologie : tout objet
n'est pas descendable. C'est la même retenue que la page 190 exprimera en
posant l'axiome \textbf{R3} --- une famille de sous-objets dont les
réalisations recouvrent est \emph{de descente effective} --- au lieu de
supposer que tout est recouvert. Le cadre le plus abstrait du carton est aussi
celui qui promet le moins, et c'est là son intérêt.

Les trois dernières lignes de la page 160 amorcent les notions d'espaces à
opérateurs admissibles et d'homomorphismes admissibles de groupes structuraux,
et s'arrêtent.

\end{document}
